\documentclass[12pt,oneside]{book}

% ============================================================
% METADATOS
% ============================================================
\newcommand{\universidad}{Universidad de Buenos Aires (UBA)}
\newcommand{\facultad}{Facultad de Derecho}
\newcommand{\carrera}{Carrera de Derecho}
\newcommand{\materia}{Derecho de la Integración}
\newcommand{\titulo}{Contratos internacionales: ley aplicable a falta de elección, vacío normativo en el Mercosur y contratación electrónica}
\newcommand{\autor}{Isaac Edgar Camacho Ocampo}
\newcommand{\anio}{2026}

% ============================================================
% PAQUETES
% ============================================================
\usepackage[utf8]{inputenc}
\usepackage[T1]{fontenc}
% Babel is intentionally omitted here because the minimal TeX installation in this environment
% does not include the Spanish language support files; the document compiles correctly with UTF-8 input alone.

\usepackage[
    top=2.5cm,
    bottom=2.5cm,
    left=3cm,
    right=2.5cm,
    headheight=15pt
]{geometry}

\usepackage{amsmath}
\usepackage{amssymb}
\usepackage{mathtools}

\usepackage{graphicx}
\usepackage{float}
\usepackage{caption}
\usepackage{subcaption}

\usepackage{booktabs}
\usepackage{array}
\usepackage{longtable}
\usepackage{tabularx}

\usepackage{enumitem}

\usepackage{xcolor}
\usepackage[most]{tcolorbox}

\usepackage{fancyhdr}
\usepackage{ragged2e}

\usepackage[
    colorlinks=true,
    linkcolor=blue!50!black,
    citecolor=green!40!black,
    urlcolor=blue!60!black,
    pdftitle={\titulo},
    pdfauthor={\autor},
    pdfsubject={Apuntes universitarios},
    bookmarks=true
]{hyperref}

% ============================================================
% CONFIGURACIÓN
% ============================================================
\graphicspath{
    {./img/}
    {./graficos/}
    {./figuras/}
}

\setcounter{tocdepth}{2}

\setlength{\parindent}{0pt}
\setlength{\parskip}{0.5em}

\setlist[itemize]{
    topsep=0.3em,
    itemsep=0.2em
}

\setlist[enumerate]{
    topsep=0.3em,
    itemsep=0.2em
}

\clubpenalty=10000
\widowpenalty=10000

\pagestyle{fancy}
\fancyhf{}
\fancyhead[L]{\nouppercase{\leftmark}}
\fancyhead[R]{\materia}
\fancyfoot[C]{\thepage}
\renewcommand{\headrulewidth}{0.4pt}
\renewcommand{\footrulewidth}{0pt}

\fancypagestyle{plain}{
    \fancyhf{}
    \fancyfoot[C]{\thepage}
    \renewcommand{\headrulewidth}{0pt}
}

% ============================================================
% COMANDOS
% ============================================================
\newcommand{\abs}[1]{\left|#1\right|}
\newcommand{\norm}[1]{\left\lVert#1\right\rVert}
\newcommand{\vect}[1]{\mathbf{#1}}
\newcommand{\deriv}[2]{\frac{d#1}{d#2}}
\newcommand{\pderiv}[2]{\frac{\partial#1}{\partial#2}}

\newcolumntype{Y}{>{\raggedright\arraybackslash}X}

% ============================================================
% ENTORNOS / CAJAS
% ============================================================
\newtcolorbox{definition}[1]{
    enhanced,
    breakable,
    title={Definición: #1},
    fonttitle=\bfseries,
    colback=blue!3,
    colframe=blue!50!black,
    boxrule=0.8pt,
    arc=2mm
}

\newtcolorbox{important}{
    enhanced,
    breakable,
    title={Importante},
    fonttitle=\bfseries,
    colback=yellow!8,
    colframe=orange!70!black,
    boxrule=0.8pt,
    arc=2mm
}

\newtcolorbox{example}[1]{
    enhanced,
    breakable,
    title={Ejemplo: #1},
    fonttitle=\bfseries,
    colback=green!4,
    colframe=green!40!black,
    boxrule=0.8pt,
    arc=2mm
}

\newtcolorbox{formula}[1]{
    enhanced,
    breakable,
    title={Fórmula / Regla: #1},
    fonttitle=\bfseries,
    colback=purple!4,
    colframe=purple!50!black,
    boxrule=0.8pt,
    arc=2mm
}

\newtcolorbox{exercise}[1]{
    enhanced,
    breakable,
    title={Ejercicio: #1},
    fonttitle=\bfseries,
    colback=gray!5,
    colframe=gray!60!black,
    boxrule=0.8pt,
    arc=2mm
}

\newtcolorbox{note}{
    enhanced,
    breakable,
    title={Nota},
    fonttitle=\bfseries,
    colback=white,
    colframe=gray!50,
    boxrule=0.6pt,
    arc=2mm
}

% ============================================================
% CONFIGURACIÓN FINAL
% ============================================================
\begin{document}

% ============================================================
% PORTADA
% ============================================================
\begin{titlepage}

\thispagestyle{empty}

\begin{center}

\vspace*{1cm}

{\large \textsc{\universidad}}

\vspace{0.2cm}

{\large \textsc{\facultad}}

\vspace{0.2cm}

{\large \carrera}

\vspace{2.5cm}

{\Huge \textbf{\materia}}

\vspace{1cm}

{\LARGE \titulo}

\vspace{0.8cm}

\textit{Comprender los conceptos es el primer paso para convertir el estudio en conocimiento.}

\vspace{1.2cm}

\rule{0.70\textwidth}{0.5pt}

\vspace{0.5cm}

{\large Apuntes de estudio}

\vspace{0.5cm}

\rule{0.70\textwidth}{0.5pt}

\vfill

{\large \autor}

\vspace{0.3cm}

{\large \anio}

\end{center}

\vfill

\begin{flushleft}
\textit{Este es un modesto aporte a los estudiantes de ``\materia'' no pretende sustituir el material oficial de la materia.}
\end{flushleft}

\end{titlepage}

% ============================================================
% ÍNDICE
% ============================================================
\tableofcontents

% ============================================================
% CONTENIDO
% ============================================================

% ------------------------------------------------------------
% CAPÍTULO 1
% ------------------------------------------------------------
\chapter[Autonomía de la voluntad]{La autonomía de la voluntad: repaso y antecedentes regionales}

El presente material continúa el estudio del contrato internacional en el Mercosur. Se parte de la autonomía de la voluntad y de sus límites, se determina qué ley rige el contrato cuando las partes no eligieron ninguna, se examina la ausencia de una norma regional común y, por último, se analiza la contratación electrónica.

\section{Repaso: autonomía material y autonomía conflictual}

El estudio del contrato internacional supone retomar tres cuestiones: el contrato, la autonomía de la voluntad y los límites a la autonomía de la voluntad. Existen la autonomía \textbf{material} y la \textbf{conflictual}, y cada una puede concebirse como \textbf{absoluta} o como \textbf{relativa}; cada una tiene su propio fundamento, y dentro de la normativa se establecen los límites.

La autonomía absoluta es aquella que, según la exposición, «viene de Dios»; la relativa es la que otorga el ordenamiento jurídico y, por ello, tiene límites. Los Estados pueden aceptar la autonomía de la voluntad con gran amplitud, poniendo pocos límites, o aceptarla poniendo muchos.

En general, los ordenamientos latinoamericanos son propensos a aceptar la autonomía material; en cambio, la autonomía conflictual se acepta de manera muy restringida o directamente no se acepta.

\section{La autonomía conflictual en el Código de Vélez y en la región}

\subsection{El Código de Vélez y la interpretación de la doctrina}

El Código de Vélez no preveía propiamente la autonomía conflictual: no establecía de modo expreso que pudiera elegirse la ley aplicable al contrato. Sin embargo, la doctrina y la jurisprudencia argentinas razonaban a partir de la autonomía material, sobre la base de que las estipulaciones insertas en los contratos son para las partes una regla a la que deben someterse como a la ley.

\begin{note}
\textbf{Referencia normativa. Código Civil de Vélez Sarsfield (derogado), art.\ 1197:} «Las convenciones hechas en los contratos forman para las partes una regla a la cual deben someterse como a la ley misma.»

En la clase no se cita el número de artículo; se transcribe el texto histórico al que remite el razonamiento. Su equivalente actual en el Código Civil y Comercial es el art.\ 959.
\end{note}

El razonamiento es el siguiente. Si lo que las partes pactan es obligatorio como la ley, podrían copiar en el contrato, por ejemplo, las obligaciones del comprador y del vendedor del Código Civil francés (``el comprador está obligado a\ldots'', ``el vendedor está obligado a\ldots'') y esas cláusulas serían obligatorias. Si las partes pueden copiar un texto normativo dentro del contrato y quedar obligadas por él como por la propia ley, \textbf{quien puede lo más, puede lo menos}: podrían incluir una sola cláusula que disponga que el contrato se rige por la ley francesa.

Así, aunque la autonomía conflictual no estuviera explícita, en la práctica la doctrina y la jurisprudencia la admitían por vía de interpretación, a partir de la posibilidad de la autonomía material.

\subsection{Paraguay y Brasil}

\begin{itemize}
    \item \textbf{Paraguay:} no interpreta la cuestión del mismo modo; la autonomía conflictual debe estar expresamente autorizada. Si no lo está, no se admite.
    \item \textbf{Brasil:} no es partidario de la autonomía conflictual; solamente admite la autonomía material.
\end{itemize}

En consecuencia, no existe una normativa específica común sobre la materia.

\section{La CIDIP de 1994 y el modelo de Roma}

En el ámbito americano existe la CIDIP, producto de la Organización de los Estados Americanos. La convención aprobada en la CIDIP de 1994 sobre derecho aplicable a los contratos tuvo solamente dos ratificaciones: la de México y la de Venezuela. México no es parte del Mercosur; solamente lo es Venezuela. Por ello, la convención no puede aplicarse en el ámbito del Mercosur.

\begin{note}
Nota editorial del apunte original: Venezuela se encuentra suspendida como Estado parte del Mercosur desde 2016.
\end{note}

La convención fue ampliamente resistida porque tomó como modelo las normas de Roma (en aquel momento, el Convenio). Preveía una autonomía bastante amplia:

\begin{itemize}
    \item las partes solamente necesitaban, en su caso, algún vínculo con la ley elegida;
    \item se permitía elegir distinta cantidad de leyes, es decir, el \textit{dépeçage};
    \item se permitía cambiar la ley aplicable en cualquier momento de la relación jurídica, incluso iniciado el juicio.
\end{itemize}

Mientras en la región se discutía si la autonomía debía admitirse o no, la CIDIP optó por copiar el modelo europeo. Ningún Estado del Mercosur la ratificó, de modo que la cuestión quedó cerrada en ese ámbito.

\begin{note}
\textbf{Referencia normativa. Convención Interamericana sobre Derecho Aplicable a los Contratos Internacionales (CIDIP V, México, 1994).}

\textbf{Art.\ 7 (primer párrafo):} «El contrato se rige por el derecho elegido por las partes.»

\textbf{Art.\ 7 (último párrafo):} «La selección de un determinado foro por las partes no entraña necesariamente la elección del derecho aplicable.»

\textbf{Art.\ 9 (primer párrafo):} «Si las partes no hubieran elegido el derecho aplicable, o si su elección resultara ineficaz, el contrato se regirá por el derecho del Estado con el cual tenga los vínculos más estrechos.»

\textbf{Art.\ 9 (segundo párrafo):} «El tribunal tomará en cuenta todos los elementos objetivos y subjetivos que se desprendan del contrato para determinar el derecho del Estado con el cual tiene vínculos más estrechos.»

Estado de ratificaciones: solo México y Venezuela, tal como lo señala la docente.
\end{note}

\begin{important}
Ni el Código de Vélez, ni Paraguay, ni Brasil, ni la CIDIP de 1994 ofrecen una regla regional común sobre la elección de la ley aplicable. La Argentina admitió la autonomía conflictual por vía de interpretación; Paraguay la exige expresa; Brasil solo admite la autonomía material; y la CIDIP, modelada sobre el Convenio de Roma, casi no tuvo ratificaciones.
\end{important}

\section{Cuadro Cornell: la autonomía de la voluntad}

\begin{longtable}{@{}
    >{\raggedright\arraybackslash}p{0.26\textwidth}
    >{\raggedright\arraybackslash}p{0.67\textwidth}@{}}
\toprule
\textbf{Preguntas / palabras clave} & \textbf{Notas / respuestas} \\
\midrule
\endfirsthead
\toprule
\textbf{Preguntas / palabras clave} & \textbf{Notas / respuestas} \\
\midrule
\endhead
\midrule
\multicolumn{2}{r}{\small\textit{continúa en la página siguiente}} \\
\endfoot
\bottomrule
\endlastfoot

\textbf{¿Qué diferencia hay entre autonomía material y autonomía conflictual?} &
En la autonomía material las partes incorporan al contrato las reglas que quieran (incluso copiando un texto normativo) y eso las obliga como la ley. En la conflictual las partes eligen directamente la ley que rige el contrato; «quien puede lo más, puede lo menos». Ambas pueden ser absolutas («vienen de Dios») o relativas (las otorga el ordenamiento y tienen límites). \\[0.6em]

\textbf{¿Cómo trataron la autonomía conflictual el Código de Vélez, Paraguay y Brasil?} &
El Código de Vélez no la preveía expresamente, pero la doctrina y la jurisprudencia argentinas la aceptaban por interpretación, a partir de la autonomía material. Paraguay exige que esté expresamente autorizada. Brasil solo admite la material. \\[0.6em]

\textbf{¿Qué fue la CIDIP de 1994 y por qué no sirve en el Mercosur?} &
Fue una convención americana (OEA) sobre derecho aplicable a los contratos, modelada sobre el Convenio de Roma: autonomía amplia, \textit{dépeçage} y cambio de ley en cualquier momento. Solo la ratificaron México y Venezuela; México no es parte del Mercosur, de modo que no se aplica en su ámbito. \\[0.6em]

\textbf{¿Existe una regla regional común sobre la elección de la ley?} &
No. Ni el Código de Vélez, ni Paraguay, ni Brasil, ni la CIDIP de 1994 la ofrecen: la Argentina la admite por interpretación, Paraguay la exige expresa, Brasil solo acepta la autonomía material y la CIDIP casi no tuvo ratificaciones. \\

\midrule
\multicolumn{2}{>{\raggedright\arraybackslash}p{\dimexpr0.93\textwidth+2\tabcolsep\relax}}{
\textbf{Resumen:} Las partes pueden incorporar reglas al contrato (autonomía material) o elegir la ley que lo rige (autonomía conflictual), absoluta o relativa según su fundamento. Los ordenamientos latinoamericanos admiten con facilidad la material y restringen la conflictual. La Argentina la admitió por interpretación del Código de Vélez; Paraguay exige autorización expresa; Brasil solo acepta la material; y la CIDIP de 1994, modelada sobre Roma, no se aplica en el Mercosur. No hay una regla regional común.
} \\
\end{longtable}

% ------------------------------------------------------------
% CAPÍTULO 2
% ------------------------------------------------------------
\chapter[Ley aplicable a falta de elección]{La ley aplicable a falta de elección de las partes}

\section{Cuándo falta la elección de ley}

Las partes pueden no haber elegido la ley aplicable por tres razones:

\begin{itemize}
    \item porque el Estado no se lo permite, o solo les permite la autonomía material y no la conflictual;
    \item porque no se les ocurrió hacerlo;
    \item porque eligieron una ley que luego resulta no aplicable.
\end{itemize}

También puede ocurrir que la elección exista pero sea \textbf{nula}: porque adolece de un defecto formal o porque, por razones de orden público, la ley elegida no es aplicable. En cualquiera de estas situaciones debe recurrirse al derecho al que envía la \textbf{norma de conflicto}; es decir, al sistema jurídico que indica conforme a qué derecho se rige una relación contractual que tiene puntos de contacto con distintos ordenamientos jurídicos, en el orden internacional.

\section{Criterio rígido o subjetivo}

Cuando los Estados fijan la ley aplicable suelen utilizar uno de dos criterios: el \textbf{rígido o subjetivo} o el \textbf{flexible u objetivo}.

\begin{definition}{Criterio rígido o subjetivo}
La ley está redactada de tal manera que indica exactamente cuál es la norma aplicable; no hay duda al respecto y el juez solo debe determinar el dato que la ley establece. Se denomina subjetivo porque se funda en una presunción sobre la voluntad de las partes: el legislador interpreta que, al elegir ese lugar, las partes se sometieron voluntariamente a ese derecho.
\end{definition}

\subsection{Los puntos de conexión del criterio rígido}

Los puntos de conexión habituales son:

\begin{itemize}
    \item \textbf{El lugar de celebración:} ``el contrato se rige por la ley del lugar de su celebración''. Brasil lo utiliza para el fondo de la cuestión y no solamente para las formas.
    \item \textbf{El lugar de cumplimiento o de ejecución:} es el criterio del derecho argentino para el fondo de la cuestión. El juez solo determina dónde se ejecuta, pero no puede variar ni aplicar otra cosa que lo que dice la ley.
\end{itemize}

En la Argentina, a la \textbf{forma} se le aplica la ley del lugar de celebración (\textit{locus regit actum}: las formas se rigen por la ley del lugar donde el acto jurídico se celebra o se realiza), mientras que al \textbf{fondo} se le aplica la ley del lugar de cumplimiento y ejecución.

\subsection{El fundamento: Savigny y la sumisión voluntaria}

Tanto para la celebración como para la ejecución el razonamiento es el mismo, pero el criterio del lugar de ejecución se adoptó siguiendo a \textbf{Savigny}. Según este autor, y conforme al principio de la \textbf{sumisión voluntaria}, debía aplicarse el derecho más conforme con la naturaleza esencial de la relación jurídica, es decir, el del ordenamiento con el cual la relación tiene más conexión: el \textbf{lugar de asiento} de la relación jurídica.

En un contrato hay dos obligaciones, y lo más importante no son las formas sino el fondo: la parte se comprometió a algo y se comprometió a cumplirlo en un determinado lugar. En ese lugar, donde se comprometió a cumplir su obligación, tiene su asiento la relación jurídica contractual. Esa conexión próxima, o contactos más próximos, se encuentra en el lugar de asiento.

No se trata de que las partes hayan dicho que consideraban esa ley, pues entonces habría autonomía de la voluntad. Las partes no dijeron nada: es el legislador quien interpreta que, al fijar que cumplirían su obligación en un lugar, se sometieron voluntariamente al derecho de ese Estado. Es una presunción lógica, un razonamiento del legislador sobre lo que la parte quería.

\begin{example}{Sumisión por el lugar de cumplimiento}
Si una parte pactó que cumpliría su obligación en Mozambique, se entiende que se sometió al derecho de Mozambique y que su obligación queda sometida a él.
\end{example}

Para la celebración el razonamiento es el mismo, aunque no se base en Savigny: si el contrato se celebró en determinado lugar, se presume que las partes se sometieron a ese derecho. Savigny hace referencia al lugar de cumplimiento porque allí está la obligación a la que la parte se compromete.

El criterio es \textbf{rígido} porque el juez, para aplicarlo, solo debe mirar lo que dice la ley y no puede modificarlo; y es \textbf{subjetivo} porque está establecido bajo la presunción de lo que las partes quisieron al contratar.

\begin{important}
Criterio rígido es igual a criterio subjetivo: la ley indica de antemano la norma aplicable (lugar de celebración o de cumplimiento) y el juez solo la aplica. Se funda en la presunción de que, al elegir ese lugar, las partes se sometieron voluntariamente a ese derecho (Savigny: sumisión voluntaria y asiento de la relación jurídica).
\end{important}

\section{Criterio flexible u objetivo}

\begin{definition}{Criterio flexible u objetivo}
No parte de la presunción de que las partes se sometieron a determinado derecho. Toma el contrato concreto, analiza todos sus elementos y el juez determina con cuál de los ordenamientos con los que está vinculado el contrato tiene el vínculo más estrecho (la conexión más próxima). El legislador dispone que el contrato se regirá por la ley del Estado con el cual presente los contactos más próximos o los vínculos más estrechos.
\end{definition}

El juez analiza, respecto de este contrato en particular y no de todos los contratos de compraventa, elementos como los siguientes:

\begin{itemize}
    \item dónde se celebró y dónde se ejecuta;
    \item dónde tienen el domicilio las partes y cuál es su nacionalidad;
    \item en qué idioma está redactado;
    \item qué figura jurídica contempla (por ejemplo, un contrato de transporte);
    \item si las partes eligieron juez o no y cuál.
\end{itemize}

Una vez establecidos los elementos y los ordenamientos con los que se vinculan, el juez determina con cuál existe el vínculo más estrecho. Va descartando lo que le parece intrascendente y reteniendo lo que le parece importante. Quien determina cuáles son los contactos más próximos y los vínculos más estrechos es el \textbf{juez}.

\begin{example}{Elementos del contrato y vínculo más estrecho}
Si ambas partes son argentinas, hay un vínculo con la Argentina; si una está en Uruguay y la otra en la Argentina, aparece Uruguay como otro elemento; si el contrato se redactó en inglés, se anota también ese vínculo. Si el contrato versa sobre una modalidad contractual regulada solamente en el derecho chino, lo más probable es que el juez concluya que el vínculo más estrecho es con China, aunque una parte sea argentina y la otra uruguaya.

En cuanto al idioma, en los contratos internacionales casi todos están redactados en inglés, con lo cual se aplicaría siempre derecho inglés; esto no siempre es así: el idioma, conjugado con otros elementos, normalmente no se toma como único elemento.
\end{example}

El criterio es \textbf{objetivo} porque no parte de una presunción sino de los elementos concretos de este contrato; y es \textbf{flexible} porque, como el análisis depende de cada contrato y del juez, este puede decidir una cosa en un caso y otra en otro, sobre una misma modalidad contractual. Así, en un contrato de compraventa el juez puede concluir que la ley aplicable es la del comprador (por el domicilio de una de las partes, el idioma de esa parte y la elección de un juez competente de ese lugar), y en otro, la del domicilio del vendedor.

% TODO: contenido incompleto o ambiguo en las notas originales (ejemplo del contrato de transporte, con datos omitidos: «desde Francia hacia […]», pasajeros «de […]», aerolínea de bandera «[…]»)
En un contrato de transporte, por ejemplo, el juez puede determinar distintas leyes según los elementos de cada caso.

\begin{table}[H]
\centering
\small
\caption{Criterio rígido (subjetivo) y criterio flexible (objetivo).}
\begin{tabularx}{\textwidth}{>{\raggedright\arraybackslash}p{0.17\textwidth} Y Y}
\toprule
\textbf{Aspecto} & \textbf{Criterio rígido o subjetivo} & \textbf{Criterio flexible u objetivo} \\
\midrule
Qué dice la ley & Indica exactamente cuál es la norma aplicable (p.\ ej., lugar de celebración o lugar de cumplimiento). & Dice que se aplica la ley del Estado con el cual el contrato presente los contactos más próximos o los vínculos más estrechos. \\
Papel del juez & Solo determina el dato (dónde se celebró o dónde se ejecuta); no puede variar ni aplicar otra cosa. & Analiza todos los elementos del contrato concreto y decide con qué ordenamiento tiene el vínculo más estrecho. \\
Fundamento & Presunción subjetiva: las partes se sometieron voluntariamente al derecho de ese lugar (Savigny). & No presume nada: parte de los elementos objetivos de este contrato (domicilios, nacionalidad, idioma, figura, juez elegido). \\
Alcance & Regla general para todos los contratos de un tipo. & Solución para este contrato en particular; puede variar entre contratos de igual modalidad. \\
Tradición jurídica & Continental romanista (países americanos). & Anglosajona, de tradición pretoriana. \\
Crítica principal & Arbitrariedad del legislador al presumir la voluntad. & Discrecionalidad del juez, posible arbitrariedad e inseguridad jurídica. \\
\bottomrule
\end{tabularx}
\end{table}

\section{Tradiciones jurídicas, críticas y estilo contractual}

El criterio rígido corresponde a la tradición \textbf{continental romanista} y el flexible a la \textbf{anglosajona}, con su tradición pretoriana, en la que son los jueces quienes elaboran y analizan.

\subsection{La crítica al criterio flexible}

Desde la tradición continental se le atribuye al sistema flexible la \textbf{discrecionalidad del juez}, que se traduce en posible arbitrariedad e \textbf{inseguridad jurídica}. Las partes pueden haber celebrado un contrato creyendo que se sometería a una regla (por ejemplo, la ley del lugar de cumplimiento) y, ante un conflicto, el juez puede concluir que los vínculos más estrechos son con otro ordenamiento (por ejemplo, la ley del domicilio del vendedor): la parte adecuó su conducta a derecho y resulta que no era el derecho aplicable.

Desde la tradición anglosajona se respondería que el criterio flexible es precisamente lo que permite establecer la ley propia de este contrato y no de todos. Esa tradición, además, descansa en una historia jurisprudencial: no se sigue lo que a un juez se le ocurre, sino que se debe fundar y elaborar muy bien la sentencia para que siente criterios jurisprudenciales. Los anglosajones confían más en la jurisprudencia y en lo que determina un juez; los sistemas continentales prefieren saber de antemano que la ley dirá, por ejemplo, ``lugar de cumplimiento'', y si no gusta, las partes eligen.

\subsection{Consecuencias en la redacción de los contratos}

\begin{itemize}
    \item \textbf{Tradición continental romanista:} se acostumbra a hacer contratos cortos, porque existen normativas específicas que resuelven los problemas. No hace falta pactar que alguien puede excusarse por incumplimiento en caso fortuito o fuerza mayor, porque cualquiera puede hacerlo si hubo ese impedimento. La ley del lugar de cumplimiento aplica caso fortuito, fuerza mayor, hecho del príncipe y las demás teorías generales, aunque no se hayan incluido en el contrato.
    \item \textbf{Tradición anglosajona:} se detalla todo. Lo que no surja del contrato se establecerá por la ley con la cual el contrato presente los vínculos más estrechos, que determinará el juez cuando ya hay conflicto; por eso conviene detallar mejor, de modo de dejar pocas cosas libradas a la incertidumbre. Se define, por ejemplo, qué es caso fortuito (el granizo, el tsunami) aunque resulte obvio. Más vale prevenir que curar.
\end{itemize}

\begin{table}[H]
\centering
\small
\caption{Tradición continental romanista y tradición anglosajona.}
\begin{tabularx}{\textwidth}{>{\raggedright\arraybackslash}p{0.19\textwidth} Y Y}
\toprule
\textbf{Aspecto} & \textbf{Continental romanista} & \textbf{Anglosajona} \\
\midrule
Criterio de ley aplicable & Rígido (subjetivo). & Flexible (objetivo). \\
Quién define la ley aplicable & El legislador, mediante una regla previa. & El juez, caso por caso, con tradición pretoriana. \\
Seguridad jurídica & Se sabe de antemano qué dirá la ley; el juez no puede apartarse. & Se confía en la jurisprudencia y en sentencias muy fundadas que sientan criterios; se critica la posible arbitrariedad. \\
Estilo de los contratos & Contratos cortos: la normativa específica y la teoría general resuelven (caso fortuito, fuerza mayor, hecho del príncipe). & Contratos largos y detallados (granizo, tsunami\ldots): lo que no se pone queda a merced de lo que decida el juez. \\
Idea que resume cada estilo & Teoría general de todo, aplicable aunque no se la haya pactado. & Más vale prevenir que curar. \\
\bottomrule
\end{tabularx}
\end{table}

Los países americanos son más bien de tradición continental romanista. La Unión Europea, en cambio, tiene la particularidad de reunir (o haber reunido) países con sistema continental romanista y países con sistema anglosajón, por lo que necesitaba aceptar los dos criterios.

% TODO: contenido incompleto o ambiguo en las notas originales (frase final inconclusa sobre Inglaterra y sus colonias: «o se centraba solamente en un criterio determinado […]»)

\section{El sistema mixto de Roma y la prestación característica}

\subsection{La técnica del modelo de Roma}

Roma combina ambos criterios en tres pasos:

\begin{enumerate}
    \item Establece una especie de \textbf{catálogo}: el contrato de compraventa se rige por determinada ley, el de locación por otra, y así sucesivamente.
    \item A falta de ello (si algún contrato no entra en ninguna de las estipulaciones, o es un híbrido que tiene parte de un contrato y parte de otro), se aplica la ley con la cual el contrato presente los \textbf{vínculos más estrechos}.
    \item Hace una \textbf{presunción}: se presume que los vínculos más estrechos están en el lugar de la \textbf{residencia habitual de la parte que tiene a su cargo el cumplimiento de la prestación característica} del contrato. Es decir, a partir de un indicio se presupone que una conexión próxima es el domicilio de quien debe cumplir la prestación más característica.
\end{enumerate}

\begin{note}
\textbf{Referencia normativa. Reglamento (CE) n.º 593/2008 «Roma I», arts.\ 3 y 4 (sucesor del Convenio de Roma de 1980).}

\textbf{Artículo 3. Libertad de elección.}

1. «El contrato se regirá por la ley elegida por las partes. Esta elección deberá manifestarse expresamente o resultar de manera inequívoca de los términos del contrato o de las circunstancias del caso. Por esta elección, las partes podrán designar la ley aplicable a la totalidad o solamente a una parte del contrato.»

2. «Las partes podrán, en cualquier momento, convenir que el contrato se rija por una ley distinta de la que lo regía con anterioridad, bien sea en virtud de una elección anterior efectuada con arreglo al presente artículo o de otras disposiciones del presente Reglamento. Toda modificación relativa a la determinación de la ley aplicable, posterior a la celebración del contrato, no obstará a la validez formal del contrato a efectos del artículo 11 y no afectará a los derechos de terceros.»

3. «Cuando todos los demás elementos pertinentes de la situación estén localizados en el momento de la elección en un país distinto de aquel cuya ley se elige, la elección de las partes no impedirá la aplicación de las disposiciones de la ley de ese otro país que no puedan excluirse mediante acuerdo.»

\textbf{Artículo 4. Ley aplicable a falta de elección.}

1. «A falta de elección realizada de conformidad con lo dispuesto en el artículo 3, y sin perjuicio de lo dispuesto en los artículos 5 a 8, la ley aplicable al contrato se determinará de este modo: a) el contrato de compraventa de mercaderías se regirá por la ley del país donde el vendedor tenga su residencia habitual; b) el contrato de prestación de servicios se regirá por la ley del país donde el prestador del servicio tenga su residencia habitual; c) el contrato que tenga por objeto un derecho real inmobiliario o el arrendamiento de un bien inmueble se regirá por la ley del país donde esté sito el bien inmueble; [\ldots] e) el contrato de franquicia se regirá por la ley del país donde el franquiciado tenga su residencia habitual; f) el contrato de distribución se regirá por la ley del país donde el distribuidor tenga su residencia habitual; g) el contrato de venta de bienes mediante subasta se regirá por la ley del país donde tenga lugar la subasta, si dicho lugar puede determinarse; [\ldots]»

2. «Cuando el contrato no esté cubierto por el apartado 1 o cuando los elementos del contrato correspondan a más de una de las letras a) a h) del apartado 1, el contrato se regirá por la ley del país donde tenga su residencia habitual la parte que deba realizar la prestación característica del contrato.»

3. «Si del conjunto de circunstancias se desprende claramente que el contrato presenta vínculos manifiestamente más estrechos con otro país distinto del indicado en los apartados 1 o 2, se aplicará la ley de este otro país.»

4. «Cuando la ley aplicable no pueda determinarse con arreglo a los apartados 1 o 2, el contrato se regirá por la ley del país con el que presente los vínculos más estrechos.»

Se transcriben los apartados más relevantes; los pasajes omitidos se indican con [\ldots]. La docente se refiere a la técnica del Convenio de Roma (catálogo, presunción de la prestación característica y vínculos más estrechos), hoy recogida en el Reglamento Roma I.
\end{note}

\subsection{Qué es la prestación más característica}

En los contratos sinalagmáticos (aquellos en que hay dos obligaciones, a diferencia de una donación) las obligaciones son de naturaleza distinta. Suele ocurrir que el pago de un precio es una de ellas; como se repite en muchos contratos, no sirve para ponerle nombre ni para categorizar el contrato. Lo que permite distinguir un contrato de otro es la prestación \textbf{no dineraria}.

\begin{definition}{Prestación más característica}
Es la prestación no dineraria, la que sirve para ponerle nombre a la relación jurídica contractual y distinguir un contrato de otro. La parte que debe cumplirla es aquella cuya residencia habitual determina la ley aplicable; el pago del precio es la prestación común y no sirve para distinguir un contrato de otro.
\end{definition}

\begin{table}[H]
\centering
\small
\caption{Prestaciones en los contratos y cuál es la «más característica».}
\begin{tabularx}{\textwidth}{>{\raggedright\arraybackslash}p{0.20\textwidth} Y Y}
\toprule
\textbf{Contrato} & \textbf{Prestación no dineraria (la característica)} & \textbf{Prestación común} \\
\midrule
Compraventa & Entrega del dominio (de la cosa). & Pago del precio. \\
Locación & Entrega del uso y goce. & Pago del precio. \\
Contrato laboral & Entrega del trabajo (obligación de hacer). & Pago del precio (la remuneración). \\
\bottomrule
\end{tabularx}
\end{table}

\subsection{¿Criterio flexible o rígido?}

Casi toda la doctrina define la teoría de la prestación más característica como una teoría dentro de un criterio flexible. Frente a esa opinión, un alumno objetó que es flexible en términos prácticos pero no deja de ser una indicación legal, como el lugar de celebración.

La posición sostenida en la exposición es que la prestación más característica es, en realidad, un criterio \textbf{rígido}, y es \textbf{el criterio más rígido de los criterios rígidos}. Si la norma dice ``residencia de la parte que tiene que cumplir la prestación más característica'', está \textbf{calificando} el lugar de cumplimiento: se entiende por lugar de cumplimiento el lugar donde el prestador más característico debe entregar la cosa, y no el del pago del precio. Entre dos obligaciones (pagar el precio o entregar la cosa) se elige una y se elimina la otra, por lo que el criterio se acota totalmente: la otra prestación no es la que interesa. No es solamente un criterio rígido que califica, por ejemplo, el domicilio de una de las partes o el lugar de cumplimiento, sino que lo acota aún más.

\section{Variantes, conveniencia y límites de la prestación característica}

\subsection{Dos variantes de la teoría}

La teoría tiene dos variantes:

\begin{enumerate}
    \item \textbf{Residencia habitual de la parte que debe cumplir la prestación más característica} (la que toma Roma).
    \item \textbf{Lugar de ejecución de la prestación más característica.}
\end{enumerate}

\begin{example}{Compraventa internacional con entrega en Uruguay}
En una compraventa internacional en la que el vendedor debe entregar la mercadería en Uruguay, el lugar de cumplimiento es donde se entrega la mercadería, no donde se paga el precio. El lugar de entrega no es el domicilio del vendedor, sino otra cosa distinta; y, sin embargo, también es prestación más característica.
\end{example}

\subsection{¿Cuál variante conviene?}

Depende. Para saber cuál de los dos criterios conviene más debe conocerse con claridad la economía del país, porque ningún país vende solamente o compra solamente:

\begin{itemize}
    \item un país \textbf{importador} se beneficia del lugar de cumplimiento de la prestación más característica, porque se le aplicará su propia ley, que conoce (la mercadería se entrega en su territorio);
    \item un país eminentemente \textbf{exportador} se beneficia de la residencia o domicilio del vendedor o prestador del servicio.
\end{itemize}

Dependerá de cuán abiertas estén las importaciones y las exportaciones para determinar qué criterio es el más conveniente; como la economía varía y cambia, no siempre es fácil, y es difícil establecer parámetros generales.

\begin{table}[H]
\centering
\small
\caption{Variantes de la prestación más característica y conveniencia económica.}
\begin{tabularx}{\textwidth}{>{\raggedright\arraybackslash}p{0.30\textwidth} Y Y}
\toprule
\textbf{Variante} & \textbf{Qué ley se aplica} & \textbf{A quién le conviene (regla general)} \\
\midrule
Residencia habitual de la parte que cumple la prestación más característica (modelo de Roma) & La del domicilio o residencia del vendedor o prestador del servicio. & País eminentemente exportador. \\
Lugar de ejecución de la prestación más característica & La del lugar donde se entrega la mercadería (p.\ ej., Uruguay). & País importador: se aplica su propia ley, que conoce. \\
\bottomrule
\end{tabularx}
\end{table}

La Unión Europea, además de los criterios generales sobre autonomía de la voluntad (artículos 3 y 4), regula específicamente cada tipo de contrato: transporte, consumo, etc. Se recomienda leer los primeros supuestos, en particular las hipótesis de prestación de servicios y de compraventa.

\subsection{El problema de la teoría: prestaciones de igual naturaleza}

La teoría de la prestación más característica sirve cuando las dos prestaciones son de distinta naturaleza. Si son de igual naturaleza, no puede aplicarse.

\begin{example}{Contratos con prestaciones de igual naturaleza}
\begin{itemize}
    \item \textbf{Depósito bancario:} se deposita dinero y se retira dinero; las dos prestaciones son dinero.
    \item \textbf{Mutuo.}
    \item \textbf{Trueque:} se cambia una cosa por otra; tampoco hay dinero.
\end{itemize}
En estos casos la prestación más característica no sirve para determinar la ley aplicable.
\end{example}

Como referencia comparada, Brasil aplica el lugar de celebración; el resto suele aplicar el lugar de cumplimiento, como los Tratados de Montevideo.

\begin{important}
A falta de elección, la ley aplicable se determina por un criterio rígido (lugar de celebración o de cumplimiento) o flexible (vínculos más estrechos). El modelo de Roma los combina: catálogo de contratos, presunción de la prestación característica y cláusula de vínculos más estrechos. La docente sostiene que la prestación característica es, en realidad, el criterio más rígido de los rígidos, y que no sirve cuando ambas prestaciones son de igual naturaleza (depósito bancario, mutuo, trueque).
\end{important}

\section{Cuadro Cornell: la ley aplicable a falta de elección}

\begin{longtable}{@{}
    >{\raggedright\arraybackslash}p{0.26\textwidth}
    >{\raggedright\arraybackslash}p{0.67\textwidth}@{}}
\toprule
\textbf{Preguntas / palabras clave} & \textbf{Notas / respuestas} \\
\midrule
\endfirsthead
\toprule
\textbf{Preguntas / palabras clave} & \textbf{Notas / respuestas} \\
\midrule
\endhead
\midrule
\multicolumn{2}{r}{\small\textit{continúa en la página siguiente}} \\
\endfoot
\bottomrule
\endlastfoot

\textbf{¿Cuándo puede faltar la elección de ley?} &
Porque el Estado solo permite autonomía material; porque no se les ocurrió a las partes; o porque eligieron una ley que resulta inaplicable (defecto formal u orden público). En todos los casos se recurre a la ley a la que remite la norma de conflicto. \\[0.6em]

\textbf{¿Qué es el criterio rígido o subjetivo?} &
La ley indica exactamente cuál es la norma aplicable (lugar de celebración o de cumplimiento) y el juez solo determina el dato. Es subjetivo porque presume que las partes se sometieron voluntariamente a ese derecho. \\[0.6em]

\textbf{¿Qué sostenía Savigny y qué siguió la Argentina?} &
Aplicar el derecho más conforme con la naturaleza esencial de la relación jurídica, es decir, el de su asiento: el lugar de cumplimiento, donde la parte se comprometió a ejecutar (sumisión voluntaria). La Argentina aplica ese lugar al fondo y la ley del lugar de celebración (\textit{locus regit actum}) a la forma; Brasil usa la celebración para ambos. \\[0.6em]

\textbf{¿Qué es el criterio flexible u objetivo?} &
No presume nada: el juez analiza todos los elementos del contrato concreto (celebración, ejecución, domicilios, nacionalidad, idioma, figura jurídica, juez elegido) y aplica la ley con la que tiene los vínculos más estrechos. Puede dar soluciones distintas para contratos de la misma modalidad. \\[0.6em]

\textbf{¿Qué críticas recibe el criterio flexible y cómo influye en los contratos?} &
Discrecionalidad del juez, posible arbitrariedad e inseguridad jurídica. Corresponde a la tradición anglosajona, que confía en la jurisprudencia: contratos largos y detallados (granizo, tsunami). La tradición continental romanista, con criterio rígido y teoría general, hace contratos cortos. \\[0.6em]

\textbf{¿Cómo combina Roma ambos criterios?} &
Catálogo de contratos con su ley; si no encajan o son híbridos, ley con los vínculos más estrechos; y una presunción: esos vínculos están en la residencia habitual de quien cumple la prestación característica. \\[0.6em]

\textbf{¿Qué es la prestación más característica y por qué la docente la considera rígida?} &
Es la no dineraria que da nombre al contrato (dominio en la compraventa, uso y goce en la locación, trabajo en el contrato laboral); el pago del precio es común y no distingue. Para la docente es el criterio más rígido de los rígidos: entre dos obligaciones elige una y elimina la otra. \\[0.6em]

\textbf{¿Residencia o lugar de ejecución? ¿A quién conviene cada una?} &
Variante 1: residencia habitual de quien cumple la prestación característica (conviene al país exportador). Variante 2: lugar de ejecución de esa prestación (conviene al país importador, que aplica su ley). Es difícil dar parámetros generales porque la economía cambia. \\[0.6em]

\textbf{¿Cuándo no sirve la prestación característica?} &
Cuando las dos prestaciones son de igual naturaleza: depósito bancario (dinero por dinero), mutuo, trueque. \\

\midrule
\multicolumn{2}{>{\raggedright\arraybackslash}p{\dimexpr0.93\textwidth+2\tabcolsep\relax}}{
\textbf{Resumen:} Cuando las partes no eligieron la ley (o su elección es inválida), se aplica la ley a la que remite la norma de conflicto, con un criterio rígido y subjetivo (celebración o cumplimiento, fundado en la sumisión voluntaria de Savigny) o flexible y objetivo (vínculos más estrechos del contrato concreto, con riesgo de arbitrariedad). Los criterios corresponden a las tradiciones continental y anglosajona y explican contratos cortos o extensos. Roma los combina con la prestación más característica, que para la docente es en el fondo un criterio muy rígido y que no funciona cuando ambas prestaciones son iguales.
} \\
\end{longtable}

% ------------------------------------------------------------
% CAPÍTULO 3
% ------------------------------------------------------------
\chapter[Síntesis y vacío en el Mercosur]{Síntesis del régimen contractual y vacío normativo en el Mercosur}

\section{Esquema de la ley aplicable al contrato}

La determinación de la ley aplicable al contrato debe resolver tres aspectos:

\begin{enumerate}
    \item Si existe la \textbf{celebración} del contrato: oferta y aceptación, y las teorías sobre el tema.
    \item La \textbf{validez formal}: si el contrato debe ser escrito o firmado; todo ello puede hacerse electrónicamente.
    \item La \textbf{validez intrínseca}, o lo que se aplica a la sustancia del contrato. Aquí se ubican la autonomía de la voluntad (material y conflictual; absoluta, que viene de Dios, o relativa, que viene del ordenamiento jurídico y por ello tiene límites) y, cuando las partes no eligieron la ley aplicable, el criterio que adopte el ordenamiento: rígido o flexible. Uno depende más de la ley; el otro, más del juez.
\end{enumerate}

\begin{table}[H]
\centering
\small
\caption{Aspectos que debe resolver la ley aplicable a un contrato internacional.}
\begin{tabularx}{\textwidth}{>{\raggedright\arraybackslash}p{0.24\textwidth} Y Y}
\toprule
\textbf{Aspecto} & \textbf{Pregunta que responde} & \textbf{Cuestiones que abarca} \\
\midrule
Existencia o celebración & ¿Se formó el contrato? & Oferta y aceptación; teorías sobre la formación del consentimiento. \\
Validez formal & ¿Se cumplieron las formas exigidas? & Si debe ser escrito, si debe ser firmado; todo ello puede hacerse electrónicamente. \\
Validez intrínseca (sustancia) & ¿Qué ley rige el fondo del contrato? & Autonomía de la voluntad (material y conflictual; absoluta o relativa, con límites). A falta de elección: criterio rígido o flexible. \\
\bottomrule
\end{tabularx}
\end{table}

\section{El vacío normativo en el Mercosur}

No existe en el Mercosur una norma específica sobre el derecho aplicable en materia contractual. Los \textbf{Tratados de Montevideo}, que sí tratan este tema, no son normativa del Mercosur ni están ratificados por todos los Estados miembros; entre la Argentina y Uruguay se aplican porque está vigente un tratado internacional, no porque sean normativa mercosureña. Venezuela, Brasil y Chile, por ejemplo, suscribieron el \textbf{Código de Bustamante} y tienen otra normativa.

\begin{note}
\textbf{Referencia normativa. Instrumentos mencionados (sin cita de artículos en clase).}

\textbf{Tratados de Montevideo de Derecho Civil Internacional:} vigentes entre la Argentina y Uruguay como tratado internacional; no constituyen normativa del Mercosur ni fueron ratificados por todos sus Estados miembros. Siguen, según la docente, el criterio del lugar de cumplimiento.

\textbf{Código de Bustamante (Código Sánchez de Bustamante):} suscripto, según la docente, por Venezuela, Brasil y Chile, que tienen por eso otra normativa.

La docente no cita artículos de estos instrumentos; por eso no se reproduce su texto.
\end{note}

La falta de una normativa clara provoca problemas en un tema que debería ser más fácil de solucionar. Si se está creando un ámbito económico, un mercado común, la negociación se realiza a través de convenios y contratos: el contrato es el instrumento que ayuda al desarrollo económico. Dejarlo sin regulación genera incertidumbre.

\begin{table}[H]
\centering
\small
\caption{Ley aplicable a los contratos según lo expuesto en clase.}
\begin{tabularx}{\textwidth}{>{\raggedright\arraybackslash}p{0.22\textwidth} Y Y}
\toprule
\textbf{Ordenamiento o instrumento} & \textbf{Autonomía de la voluntad} & \textbf{Ley aplicable a falta de elección} \\
\midrule
Argentina (precedente: Código de Vélez) & Material expresa; conflictual aceptada por interpretación de la doctrina y la jurisprudencia. & Fondo: ley del lugar de cumplimiento o ejecución (Savigny). Forma: ley del lugar de celebración (\textit{locus regit actum}). \\
Paraguay & La conflictual debe estar expresamente autorizada. & --- \\
Brasil & Solo autonomía material; no es partidario de la conflictual. & Ley del lugar de celebración, tanto para el fondo como para la forma. \\
CIDIP de 1994 (México) & Amplia, modelada sobre Roma. Solo la ratificaron México y Venezuela. & --- \\
Convenio de Roma / Unión Europea & Amplia: \textit{dépeçage}, cambio de ley en cualquier momento, vínculo mínimo. & Catálogo de contratos, presunción de la prestación característica y vínculos más estrechos. \\
Tratados de Montevideo & Vigentes entre Argentina y Uruguay como tratado, no como norma del Mercosur. & Lugar de cumplimiento. \\
Código de Bustamante & Suscripto por Venezuela, Brasil y Chile. & --- \\
\bottomrule
\end{tabularx}
\end{table}

\begin{important}
El Mercosur carece de una norma propia sobre ley aplicable a los contratos. Para un mercado común, donde el contrato es el instrumento del desarrollo económico, esa falta de regulación genera incertidumbre.
\end{important}

\section{Cuadro Cornell: síntesis y vacío en el Mercosur}

\begin{longtable}{@{}
    >{\raggedright\arraybackslash}p{0.26\textwidth}
    >{\raggedright\arraybackslash}p{0.67\textwidth}@{}}
\toprule
\textbf{Preguntas / palabras clave} & \textbf{Notas / respuestas} \\
\midrule
\endfirsthead
\toprule
\textbf{Preguntas / palabras clave} & \textbf{Notas / respuestas} \\
\midrule
\endhead
\midrule
\multicolumn{2}{r}{\small\textit{continúa en la página siguiente}} \\
\endfoot
\bottomrule
\endlastfoot

\textbf{¿Qué aspectos rige la ley aplicable al contrato?} &
Existencia o celebración (oferta y aceptación), validez formal (si debe ser escrito o firmado, lo cual puede hacerse electrónicamente) y validez intrínseca o sustancia del contrato (autonomía de la voluntad y, a falta de elección, criterio rígido o flexible). \\[0.6em]

\textbf{¿Cuál es la situación del Mercosur en materia de ley aplicable a los contratos?} &
No tiene una norma específica. Los Tratados de Montevideo rigen entre Argentina y Uruguay como tratado internacional, no como norma mercosureña, y no están ratificados por todos los Estados miembros; Venezuela, Brasil y Chile suscribieron el Código de Bustamante, que es otra normativa. \\[0.6em]

\textbf{¿Qué diferencias de criterio hay entre los ordenamientos comparados?} &
A falta de elección: la Argentina aplica el lugar de cumplimiento al fondo y el de celebración a la forma; Brasil, el de celebración para ambos; los Tratados de Montevideo, el lugar de cumplimiento; Roma combina catálogo, prestación característica y vínculos más estrechos. En cuanto a la autonomía, Paraguay exige autorización expresa y Brasil solo admite la material. \\[0.6em]

\textbf{¿Por qué la falta de una norma regional común complica la integración?} &
Porque en un mercado común se negocia a través de convenios y contratos, que son el instrumento del desarrollo económico; si no se los regula, hay incertidumbre en un tema que debería ser más fácil de solucionar. \\

\midrule
\multicolumn{2}{>{\raggedright\arraybackslash}p{\dimexpr0.93\textwidth+2\tabcolsep\relax}}{
\textbf{Resumen:} La ley aplicable al contrato debe resolver la existencia, la validez formal y la validez intrínseca. Sobre esos aspectos los ordenamientos de la región adoptan soluciones dispares y el Mercosur carece de una norma propia: solo hay tratados dispersos (Montevideo, Bustamante) que no son normativa del bloque. Para un mercado común, donde el contrato es el instrumento del desarrollo económico, esa ausencia debilita la seguridad.
} \\
\end{longtable}

% ------------------------------------------------------------
% CAPÍTULO 4
% ------------------------------------------------------------
\chapter[Contratación electrónica]{Contratación electrónica}

\section{Naturaleza jurídica de internet}

Al surgir el comercio electrónico y la contratación electrónica, el primer interrogante fue cuál es la naturaleza de internet. La doctrina se dividió desde el comienzo en dos posturas.

\subsection{Internet como espacio}

Algunos entendían que internet era un \textbf{espacio}. La terminología empleada lo sugiere: se habla de domicilio electrónico, de enviar un mail como se enviaría una carta a una dirección, de la plataforma como si fuera un lugar físico, del ``sitio'' de internet. Se sostenía entonces que era un espacio totalmente nuevo, independiente del territorio nacional, que no pertenecía al territorio de ningún Estado.

La consecuencia lógica era que, si no formaba parte del territorio, no se le podía aplicar el derecho estatal; los propios operadores de internet se regirían por sus propias normas, con una amplia autonomía de la voluntad de las partes: lo que las partes establecieran se haría. Esta idea parecía ideal, pero tenía problemas:

\begin{itemize}
    \item a través de internet no solamente se compra y se vende: también se cometen delitos, como pornografía infantil y difamación, y es difícil que por autonomía de la voluntad se llegue a un acuerdo sobre las reglas;
    \item no solo no había derecho, sino que tampoco se sabía qué jueces eran competentes para resolver los conflictos que se dan en el ámbito de internet.
\end{itemize}

Esta corriente ya no es tan fervorosa como antes, porque se advierte que, aunque la teoría funcione, en la práctica se necesitan ciertas normas y acudir a un juez.

\subsection{Internet como medio de comunicación}

La otra postura sostiene que internet no es un territorio ni un sitio: es un \textbf{medio de comunicación} como cualquier otro. En distintos momentos las personas se comunicaron mediante señales de humo, paloma mensajera, cartas, telégrafo, teléfono y fax; en todos los casos la voluntad de una parte llegaba a la otra.

Si es un medio de comunicación y no un sitio, así como se regula la carta (que debe estar firmada), también puede regularse este medio y exigirle lo que debe garantizar para que sirva para comunicarse.

\begin{table}[H]
\centering
\small
\caption{Dos posturas sobre la naturaleza de internet.}
\begin{tabularx}{\textwidth}{>{\raggedright\arraybackslash}p{0.20\textwidth} Y Y}
\toprule
\textbf{Aspecto} & \textbf{Internet como espacio} & \textbf{Internet como medio de comunicación} \\
\midrule
Idea central & Un sitio nuevo, anacional, independiente del territorio de los Estados. & Un medio más, como las señales de humo, la paloma mensajera, la carta, el telégrafo, el teléfono o el fax. \\
Derecho aplicable & No se aplica el derecho estatal; los operadores se rigen por sus propias normas (amplia autonomía de la voluntad). & Se aplica el ordenamiento jurídico; se le exige al medio que garantice los requisitos legales. \\
Jueces competentes & No se sabía qué jueces eran competentes. & Los que corresponden según las reglas generales. \\
Problemas & En la práctica se necesitan normas y acudir a un juez: delitos, difamación, pornografía infantil. & --- \\
\bottomrule
\end{tabularx}
\end{table}

\section{Inexistencia de un derecho informático autónomo}

Según la posición expuesta, no existe un derecho informático como derecho situado fuera del sistema jurídico y aplicable a todos. La informática atraviesa el derecho de manera \textbf{transversal}: pueden cometerse delitos, celebrarse contratos, pagarse la cuota alimentaria o contraerse matrimonio a través de un medio informático. Es un medio de comunicación que permite realizar las mismas actividades que regula el ordenamiento jurídico.

La doctrina sostenía que había que crear reglas nuevas porque el fenómeno era novedoso. La cuestión es si se necesitan reglas nuevas o si debe verificarse que el medio de comunicación permita cumplir el requisito que la ley establece para la validez del acto. Se sostiene lo segundo: \textbf{leyes nuevas, no}; se le debe exigir al medio que garantice determinadas cuestiones, por ejemplo la identidad del emisor, de modo que quien recibe un mensaje sepa que proviene de quien dice ser.

En tanto el medio electrónico se actualice y brinde la seguridad que requiere el ámbito jurídico, es un medio utilizado por el derecho, como lo fueron el fax, la carta, el teléfono o el telégrafo (con puntos y rayas).

Existe abundante normativa elaborada por las Naciones Unidas, a través de la Comisión de las Naciones Unidas para el Derecho Mercantil Internacional, sobre documento electrónico y firma digital, pero todas se reducen a tres o cuatro principios generales, que resuelven la cuestión.

\begin{note}
\textbf{Referencia normativa. Normativa de la CNUDMI (Naciones Unidas).} La docente alude, sin citar artículos, a las normas de la Comisión de las Naciones Unidas para el Derecho Mercantil Internacional sobre documento electrónico y firma digital, y afirma que todas se reducen a tres o cuatro principios: neutralidad tecnológica, no discriminación y equivalencia funcional.
\end{note}

\begin{important}
No hace falta un «derecho informático» autónomo: la informática atraviesa transversalmente todas las ramas del derecho. En lugar de crear leyes nuevas, se le exige al medio electrónico que garantice lo que el derecho requiere (identidad del emisor, acceso a la manifestación de voluntad, integridad).
\end{important}

\section{Neutralidad tecnológica y no discriminación}

\begin{definition}{Principio de neutralidad tecnológica}
Al regular el medio informático, así como se reguló la carta (con el sello postal, siendo el correo un medio de comunicación), no debe hacerse pensando en una programación específica. La tecnología cambia tan rápido que, si se tiene en vista un determinado sistema o modelo de documento, la regulación queda superada antes de terminarse.
\end{definition}

\begin{example}{El formato PDF y el documento original}
Cuando surgió internet, los archivos PDF no podían modificarse. Por ello algunos propusieron que, cuando la ley exigiera documento original, debía presentarse el PDF, porque no podía modificarse ni falsificarse. La versión siguiente de Windows permitió modificarlo, de modo que ya no puede afirmarse que el PDF equivalga al documento original. Esto muestra que la tecnología cambia tan rápido que los requisitos o las medidas no pueden tomarse en base a un sistema determinado.
\end{example}

\begin{definition}{Principio de no discriminación}
No se debe discriminar entre la utilización de un medio electrónico y un medio tradicional. Si se los solicita a los fines de prueba, tienen tanto valor uno como otro; y si se los solicita a los fines de la realización del acto, tiene tanto valor que la voluntad se manifieste por un medio o por el otro. La no discriminación no es solamente a los fines de la prueba, sino también a los fines de la propia existencia del acto.
\end{definition}

Si en materia de oferta y aceptación el consentimiento se envía a través de un mail, no para la prueba sino para que el acto exista, se rige por las mismas reglas: la ley no discrimina.

Quedaba por resolver cómo utilizar un medio electrónico si es inseguro y no garantiza múltiples cuestiones. La respuesta está en la matriz de estos principios: la equivalencia funcional.

\section{Equivalencia funcional: escrito y original}

\begin{definition}{Principio de equivalencia funcional}
Cuando el derecho exige un requisito para realizar un acto, no lo exige porque sí, sino porque busca una finalidad específica. Si un medio electrónico permite cumplir esa finalidad, es el equivalente funcional del medio tradicional. No es un derecho creado independientemente: es un medio que se utiliza en el ámbito del derecho.
\end{definition}

\subsection{El requisito de «por escrito»}

La ley exige que algo conste por escrito, en contraposición con la exteriorización verbal, para tener constancia y poder acceder nuevamente a la manifestación de voluntad exactamente en los términos en que se hizo. Se exige para la existencia y la validez del acto: si no se cumple, el acto es inválido. Escrito se contrapone a la manifestación verbal: la palabra se la lleva el viento. Si algo se puso por escrito, es exactamente lo que está por escrito, sin interpretación alguna.

Por tanto, cualquier medio tecnológico que permita volver a acceder a la manifestación de voluntad que hizo la parte, exactamente como la hizo, es el equivalente funcional del medio escrito (el celular, la computadora, cualquier medio electrónico). No hace falta imprimirlo. Un mail es escrito no porque pueda imprimirse, sino porque se puede abrir la casilla de entrada o la de salida y volver a leerlo.

\subsection{El requisito de «original»}

La exigencia de original no se satisface con el mero hecho de que algo conste por escrito: se solicita el original para poder constatar que nadie haya adulterado el documento. Así, en un escrito oficial se podía borrar raspando con una cuchilla y poner corrector, y eso se advertía al ver el original.

Si un medio electrónico permite verificar si existieron o no adulteraciones, o si es posible hacerlas, el documento es original. El \textbf{documento original} es el que se emitió por primera vez, tal cual los términos en que fue emitido y que guarda la computadora.

\begin{itemize}
    \item El mail en la bandeja de \textbf{salida} es el documento original del emisor.
    \item El mail en la bandeja de \textbf{entrada} es el otro ejemplar, en un mismo tenor y a un mismo efecto; ambos son originales.
    \item Si el destinatario responde en el mismo mail con una contraoferta, la respuesta es su documento original (su contrapropuesta); el texto anterior que aparece al pie ya no es documento original, es una copia, porque no es el que se recibió por primera vez sino un reenvío y puede borrarse o modificarse.
    \item El mail recibido en la bandeja de entrada no puede modificarse: para modificarlo hay que reenviárselo.
\end{itemize}

El medio electrónico permite, entonces, tener documentos que se envían y reciben y no son susceptibles de modificación: eso es original. Todo lo demás son copias.

\begin{example}{Error en el monto de un contrato enviado por mail}
Un alumno plantea el caso de una persona que acuerda un contrato, lo prepara con inteligencia artificial, comete un error en el monto (acordaron 500.000 y se agrega un cero de más) y lo envía igual. La respuesta fue que, si se lo envió, es válido, porque es su manifestación de voluntad: al enviar el mail se hizo la manifestación de voluntad. Si se trata de un error de tipeo o error material y se lo advierte y se lo subsana de inmediato, entra una cuestión distinta (véase la rectificación en la sección sobre momento y lugar).
\end{example}

\section{Equivalencia funcional: la firma}

La firma se exige para corroborar que la persona es la que manifestó la voluntad y para vincular a una persona determinada con una manifestación de voluntad. Si el texto dice que «Juan Pérez dijo tal cosa», la firma sirve para identificar al que manifestó la voluntad y relacionarla con el texto. La firma tiene, por tanto, un doble efecto: \textbf{identificar a la parte} y \textbf{vincularla con una manifestación de voluntad}.

\begin{definition}{Firma (equivalente funcional electrónico)}
Cualquier medio electrónico que permita identificar a un sujeto, vincularlo con el contenido de un mensaje y que esté bajo su \textbf{exclusivo control}, según dice la ley, es firma, porque cumple la finalidad del requisito.
\end{definition}

La mano de la persona está bajo su exclusivo control: nadie puede firmar por ella. Del mismo modo, la firma digital (ese código, esa forma de firmar) está bajo el control del firmante. Si el firmante entrega su clave a otro y este la usa, la consecuencia recae sobre el firmante.

No se están creando reglas nuevas para la informática: se toman los mismos actos jurídicos y se parte de que cada requisito se exige con determinada finalidad. Si esa finalidad se cubre con el medio informático, este es el equivalente funcional del medio tradicional; no hace falta una ley nueva que establezca que la firma electrónica es válida, porque puede cumplirse el requisito.

\begin{table}[H]
\centering
\small
\caption{Principio de equivalencia funcional: requisito, finalidad y equivalente electrónico.}
\begin{tabularx}{\textwidth}{>{\raggedright\arraybackslash}p{0.14\textwidth} Y Y}
\toprule
\textbf{Requisito legal} & \textbf{Finalidad que persigue la ley} & \textbf{Equivalente funcional electrónico} \\
\midrule
Por escrito & Tener constancia y poder acceder de nuevo, exactamente, a los términos en que se manifestó la voluntad («la palabra se la lleva el viento»). & Cualquier medio (celular, computadora) que permita volver a acceder a la manifestación tal como se hizo. Un mail es escrito porque se lo puede releer en la casilla de entrada o de salida; no hace falta imprimirlo. \\
Original & Constatar que nadie adulteró el documento. & El documento emitido por primera vez y que guarda la computadora: el mail de la bandeja de salida es el original; el de la bandeja de entrada es el otro ejemplar. El texto citado al pie de una respuesta es copia. \\
Firma & Identificar a quien manifestó la voluntad y vincularlo con el texto. & Cualquier medio que identifique a un sujeto, lo vincule con el contenido del mensaje y esté bajo su exclusivo control. \\
\bottomrule
\end{tabularx}
\end{table}

\begin{important}
Tres principios resuelven casi todo: neutralidad tecnológica (no regular pensando en un sistema específico), no discriminación (mismo valor para la prueba y para la existencia del acto) y equivalencia funcional (si el medio cumple la finalidad por la que la ley exige el requisito, es el equivalente funcional del medio tradicional).
\end{important}

\section{Firma digital y firma electrónica en la Argentina y en el derecho comparado}

\subsection{El régimen argentino}

Si existiera una ley que receptara los principios, se solucionarían casi todos los problemas del uso de tecnologías, y no habría que modificarla cada vez que cambia el PDF o el sistema.

La Argentina distingue entre \textbf{firma digital} y \textbf{firma electrónica}:

\begin{itemize}
    \item La firma digital goza de \textbf{presunción de validez}.
    \item En la firma electrónica ocurre lo contrario: quien afirma que esa firma le corresponde a alguien debe probarlo.
\end{itemize}

Al legislar sobre la firma digital, la Argentina \textbf{no respetó el principio de neutralidad tecnológica}. Estableció que por firma digital se entiende el sistema de criptografía asimétrica de cruce de clave pública con clave privada, certificado por un ente certificador. Ese sistema de identificación es, para el ordenamiento argentino, el que tiene exacta equivalencia funcional con la firma manuscrita; es un sistema específico. La clave privada es la que se activa al colocar el token; la clave pública es la que dio la empresa que maneja el sistema. Ese sistema es el único que se equipara a la firma manuscrita; respecto de cualquier otro, debe probarse.

\begin{note}
\textbf{Referencia normativa. Ley 25.506 de Firma Digital (Argentina), arts.\ 2, 3, 5 y 6 (texto original).}

\textbf{Art.\ 2. Firma digital:} «Se entiende por firma digital al resultado de aplicar a un documento digital un procedimiento matemático que requiere información de exclusivo conocimiento del firmante, encontrándose ésta bajo su absoluto control. La firma digital debe ser susceptible de verificación por terceras partes, tal que dicha verificación simultáneamente permita identificar al firmante y detectar cualquier alteración del documento digital posterior a su firma.»

\textbf{Art.\ 3. Del requerimiento de firma:} «Cuando la ley requiera una firma manuscrita, esa exigencia también queda satisfecha por una firma digital. Este principio es aplicable a los casos en que la ley establece la obligación de firmar o prescribe consecuencias para su ausencia.»

\textbf{Art.\ 5. Firma electrónica:} «Se entiende por firma electrónica al conjunto de datos electrónicos integrados, ligados o asociados de manera lógica a otros datos electrónicos, utilizado por el signatario como su medio de identificación, que carezca de alguno de los requisitos legales para ser considerada firma digital. En caso de ser desconocida la firma electrónica corresponde a quien la invoca acreditar su validez.»

\textbf{Art.\ 6. Documento digital:} «Se entiende por documento digital a la representación digital de actos o hechos, con independencia del soporte utilizado para su fijación, almacenamiento o archivo. Un documento digital también satisface el requerimiento de escritura.»

Texto original de la Ley 25.506; algunos de sus artículos fueron modificados por la Ley 27.446 (2018). La docente no cita números de artículo; se transcriben los que definen los conceptos que explica.
\end{note}

\begin{table}[H]
\centering
\small
\caption{Firma digital y firma electrónica en la Argentina.}
\begin{tabularx}{\textwidth}{>{\raggedright\arraybackslash}p{0.22\textwidth} Y Y}
\toprule
\textbf{Aspecto} & \textbf{Firma digital} & \textbf{Firma electrónica} \\
\midrule
Sistema & Criptografía asimétrica: clave privada (p.\ ej., el token) y clave pública, certificada por un ente certificador. & Cualquier otro medio de identificación electrónica que no cumpla los requisitos de la firma digital. \\
Presunción & Goza de presunción de validez. & Al revés: quien invoca la firma debe probar que le corresponde a esa persona. \\
Equiparación con la firma manuscrita & Es el único sistema que la Argentina equipara a la firma manuscrita. & Hay que probar su validez. \\
Crítica de la docente & No respeta la neutralidad tecnológica: fija un sistema técnico específico. & --- \\
\bottomrule
\end{tabularx}
\end{table}

\subsection{El criterio de los ordenamientos modelo y el caso de Chile}

Tanto la ley modelo como las directivas de la Unión Europea que regulan el comercio y la contratación electrónicos, así como algunos ordenamientos nacionales que se basaron en las normas internacionales, adoptan un criterio que no toma un sistema específico, sino aquello que permita identificar a un sujeto, vincularlo con un texto y que esté bajo su exclusivo control.

La ley de \textbf{Chile} toma esos parámetros y establece que \textbf{audio, voz o video} constituyen firma. Un mensaje de voz enviado por celular contiene la voz de la persona con los términos exactos de lo que dijo; puede recuperarse cuantas veces se quiera exactamente lo dicho, por lo que sería funcionalmente equivalente a escrito, y quien lo invoca puede probar que es esa voz. A la jurisprudencia argentina le costará comprender esta solución; pero corresponde que quien lo niegue pruebe que no es su voz.

\begin{note}
\textbf{Referencia normativa. Ley 19.799 de Chile sobre documentos electrónicos y firma electrónica, art.\ 2 f) y art.\ 3 (último inciso).}

\textbf{Art.\ 2, letra f):} «Firma electrónica: cualquier sonido, símbolo o proceso electrónico, que permite al receptor de un documento electrónico identificar al menos formalmente a su autor;»

\textbf{Art.\ 3, último inciso:} «La firma electrónica, cualquiera sea su naturaleza, se mirará como firma manuscrita para todos los efectos legales, sin perjuicio de lo establecido en los artículos siguientes.»

La docente se refiere a la ley chilena sin citar artículos y resume su criterio como «audio, voz o video»; se transcriben las disposiciones que consagran ese criterio amplio.
\end{note}

\subsection{Video, presencia y reconocimiento facial}

En la práctica, las audiencias penales y otras se toman de manera virtual. Cuando una persona participa en una audiencia virtual, se le exige que se ubique frente a la cámara, porque es necesario verla para saber que es ella quien manifiesta: es el equivalente funcional. Aunque no esté establecido específicamente, resulta correcto si se aplican los principios generales.

\subsection{Hacia la presencialidad: la conexión ininterrumpida}

La exigencia de presencialidad se vincula con que la conexión sea ininterrumpida o no. En una videoconferencia, salvo una interrupción de milésimas de segundo, se entiende que la conexión es ininterrumpida; por lo tanto, las manifestaciones se hacen entre presentes. En cambio, si hay una interrupción en la comunicación (un mail sale de un sistema y entra a otro distinto, aunque pueda hacerlo en un instante), eso es \textbf{transporte}.

\section{Momento, lugar, directivas europeas, Mercosur y cierre}

Para entender la contratación o los negocios electrónicos no debe pensarse que haya que crear todos los ordenamientos jurídicos nuevos: debe preguntarse qué requisito exige la ley y para qué lo exige. Si se lo puede cumplir con un medio electrónico, este es el equivalente funcional que sirve.

\subsection{La rectificación}

Si quien envió un mail advierte de inmediato que cometió un error, puede enviar la corrección (``perdón, me equivoqué''). Es una \textbf{rectificatoria}, prevista dentro de lo normativo y que resulta posible: al haberse rectificado de inmediato, la rectificación es válida.

\subsection{Qué garantizan los tres principios}

Neutralidad tecnológica, equivalencia funcional y no discriminación son los tres principios principales; para algún caso particular podría mencionarse algún otro. Garantizan, por ejemplo:

\begin{itemize}
    \item Con la neutralidad tecnológica no puede negarse la validez de un audio ni decirse ``yo no lo dije'': es lo mismo que si se hubiera firmado. Así como ante una firma se llama a un perito calígrafo para verificarla, aquí se llamará a un perito tecnológico para determinar si se adulteró la voz.
    \item Con la no discriminación no se puede discriminar ni a los fines de la existencia ni a los fines de la prueba.
    \item Con la equivalencia funcional, si se logra cumplir la finalidad que tenía la ley al establecer el requisito, es el equivalente funcional de ese requisito (escrito, original, firma).
\end{itemize}

\subsection{Las directivas de la Unión Europea: momentos y lugares}

En la Unión Europea existen numerosas directivas que regulan estas cuestiones, todas bajo el principio de equivalencia funcional. No cambian el criterio, sino que \textbf{precisan momentos y lugares}: cuándo se entiende que salió el mensaje y cuándo se interpreta que llegó.

\begin{itemize}
    \item El mensaje \textbf{salió} cuando se hizo ``enter'' y salió del control de quien lo envió: la manifestación se efectuó.
    \item El mensaje \textbf{llegó} cuando entró en un sistema que está bajo el control del destinatario (por ejemplo, el servidor que está a su disposición o el servidor ubicado en otro lugar, como Nueva York); figura la fecha de entrada.
\end{itemize}

\subsection{Audio, voz, video y reconocimiento facial}

También se habla de audio, voz y video, de allí la utilización de la huella digital en una pantalla, el iris del ojo y el reconocimiento facial. Es el caso de la verificación facial que se solicita en ARCA: el sistema identifica a la persona. Lo que identifica a la persona (su firma, su voz) cumple la función de la firma.

\subsection{El Mercosur}

El Mercosur tiene algunas disposiciones que toman la firma electrónica y la firma digital como equivalentes y les dan plena validez, pero no profundizó mucho más: no estableció una autoridad de certificación única.

\begin{note}
\textbf{Referencia normativa. Disposiciones del Mercosur sobre firma electrónica y digital.} La docente señala, sin citar normas ni artículos, que el Mercosur tiene algunas disposiciones que toman la firma electrónica y la firma digital como equivalentes y les reconocen plena validez, sin haber establecido una autoridad de certificación única.
\end{note}

\begin{important}
Los contratos internacionales plantean tres preguntas encadenadas: (1) ¿pueden las partes elegir la ley? (autonomía material y conflictual, con límites); (2) si no la eligieron, ¿qué ley se aplica? (criterio rígido o flexible; modelo mixto de Roma con prestación característica); y (3) ¿vale un contrato celebrado por medios electrónicos? (sí, por neutralidad tecnológica, no discriminación y equivalencia funcional). El Mercosur no tiene una norma propia para ninguna de las dos primeras.
\end{important}

\section{Cuadro Cornell: contratación electrónica}

\begin{longtable}{@{}
    >{\raggedright\arraybackslash}p{0.26\textwidth}
    >{\raggedright\arraybackslash}p{0.67\textwidth}@{}}
\toprule
\textbf{Preguntas / palabras clave} & \textbf{Notas / respuestas} \\
\midrule
\endfirsthead
\toprule
\textbf{Preguntas / palabras clave} & \textbf{Notas / respuestas} \\
\midrule
\endhead
\midrule
\multicolumn{2}{r}{\small\textit{continúa en la página siguiente}} \\
\endfoot
\bottomrule
\endlastfoot

\textbf{¿Qué es internet para el derecho?} &
Hay dos posturas: un espacio nuevo, anacional, que se rige por la autonomía de la voluntad (hoy menos fervorosa porque en la práctica hacen falta normas y jueces), o un medio de comunicación como la carta, el telégrafo o el fax, que el derecho regula y al que le exige garantías. \\[0.6em]

\textbf{¿Hace falta un derecho informático autónomo?} &
No. La informática atraviesa transversalmente el derecho (delitos, contratos, alimentos, matrimonio). No se crean leyes nuevas: se le exige al medio que garantice la identidad del emisor y el cumplimiento de los requisitos legales. \\[0.6em]

\textbf{¿Cuáles son los tres principios y qué significa cada uno?} &
Neutralidad tecnológica: no regular pensando en un sistema específico (ejemplo del PDF). No discriminación: igual valor del medio electrónico para la prueba y para la existencia del acto. Equivalencia funcional: si el medio cumple la finalidad del requisito, es su equivalente. \\[0.6em]

\textbf{¿Cuándo un mail es «escrito» y «original»?} &
Es escrito si permite volver a acceder exactamente a la manifestación de voluntad (no hace falta imprimirlo). Original es el documento emitido por primera vez, que no puede ser adulterado: el de la bandeja de salida (y el de la de entrada, como el otro ejemplar); el texto al pie de una respuesta es copia. Un mail enviado con error de monto es manifestación de voluntad; solo se subsana rectificando de inmediato. \\[0.6em]

\textbf{¿Qué es la firma y cuándo hay firma digital o electrónica?} &
La firma identifica al autor y lo vincula con el texto; cualquier medio que lo logre y esté bajo su exclusivo control es firma. La Argentina distingue firma digital (criptografía asimétrica certificada; presunción de validez) de firma electrónica (se debe probar) y no respetó la neutralidad tecnológica. La ley chilena admite audio, voz o video. \\[0.6em]

\textbf{¿Cuándo salió y cuándo llegó un mensaje electrónico? ¿Qué dice el Mercosur?} &
Salió cuando sale del control del emisor (el ``enter''); llegó cuando ingresa a un sistema bajo control del destinatario. Las directivas europeas precisan esos momentos y lugares. El Mercosur equipara firma electrónica y digital y les da validez, sin una autoridad certificante única. \\

\midrule
\multicolumn{2}{>{\raggedright\arraybackslash}p{\dimexpr0.93\textwidth+2\tabcolsep\relax}}{
\textbf{Resumen:} En la contratación electrónica no hacen falta leyes nuevas ni un derecho informático autónomo: internet es un medio de comunicación al que se le exige cumplir las finalidades de los requisitos legales, y tres principios lo resuelven (neutralidad tecnológica, no discriminación y equivalencia funcional), de modo que un mail puede ser escrito y original, y un medio que identifique al autor bajo su exclusivo control puede ser firma. La Argentina distingue firma digital y electrónica con un sistema técnico específico; Chile admite audio, voz y video; la Unión Europea precisa momentos y lugares; y el Mercosur apenas equiparó ambas firmas.
} \\
\end{longtable}

\end{document}
